\title{FlashAttention-3: Fast and Accurate Attention with Asynchrony and Low-precision}

\begin{abstract}
Attention remains a fundamental component of the ubiquitous Transformer framework, yet it often constitutes a performance bottleneck for large language models and long-context processing tasks. \textsc{FlashAttention} introduced a GPU-optimized strategy for accelerating attention by reducing memory accesses; however, it does not leverage the advancements found in the latest GPU hardware, with \textsc{FlashAttention-2} utilizing only 35\% of the available capacity on the H100 GPU. In this work, we introduce three novel approaches to further accelerate attention on Hopper GPUs: utilizing Tensor Core and TMA asynchrony to (1) enable warp-specialized computation that overlaps with data transfer, (2) intermix block-level matrix multiplication and softmax steps, and (3) perform block-wise quantization and incoherent computation, capitalizing on hardware support for FP8 low-precision operations. The resulting approach, \textsc{FlashAttention-3}, delivers a 1.5-2.0$\times$ increase in throughput on H100 GPUs, with FP16 workloads achieving up to 740 TFLOPs/s (corresponding to 75\% utilization), and FP8 workloads reaching nearly 1.2 PFLOPs/s. Additionally, our experiments show that FP8 \textsc{FlashAttention-3} realizes a 2.6$\times$ reduction in numerical error compared with a standard FP8 attention implementation.
\end{abstract}

\section{Introduction}
\label{sec:intro}

Within the Transformer architecture~\citep{vaswani2017attention}, the attention module represents the main source of computational expense, primarily due to the self-attention operation between queries and keys exhibiting quadratic complexity with respect to sequence length. Enabling attention over extended contexts can provide new functionalities, such as improved modeling and reasoning across multiple lengthy documents~\citep{guo2021longt5,shaham2022scrolls,peng2023yarn}, analysis of extensive code repositories~\citep{roziere2023code, li2023starcoder}, as well as accommodating diverse modalities like high-resolution visual data~\citep{chen2022scaling}, audio signals~\citep{gulati2020conformer}, and video streams~\citep{ho2022video}. Additionally, it opens possibilities for applications demanding the processing of extended sequential histories, such as user interactions~\citep{sun2019bert4rec} and agents planning over long time horizons~\citep{yao2022react}. These opportunities have motivated substantial research into accelerating attention computations in long-sequence scenarios, pursued via approximation techniques~\citep{katharopoulos2020transformers,choromanski2020rethinking,
tay2020efficient}, efficient software solutions (\citep{rabe2021self,dao2022flashattention,kwon2023efficient}), and the exploration of novel architectural paradigms~\citep{peng2023rwkv,sun2023retentive,gu2023mamba}.

This paper extends the contributions of \citet{dao2022flashattention}, who proposed exact-attention algorithms leveraging detailed insights into the GPU execution model and associated hardware properties when designing efficient approaches. In their seminal work \citep{dao2022flashattention}, Dao et al. presented \textsc{FlashAttention}, which introduced a novel tiling method enabling all attention computations to be fused within a single GPU kernel, thereby avoiding the costly process of reading and writing to global memory between intermediate steps. 

Building on these concepts, \citet{dao2023flashattention2} redesigned the strategy as \textsc{FlashAttention-2}, which further extends parallelization to the sequence length and organizes the forward pass’s inner loop to operate across blocks of the key and value matrices. These enhancements deliver improved utilization and more effective workload allocation on GPU hardware. Nevertheless, our analysis indicates that, on modern GPUs, \textsc{FlashAttention-2} still falls short in resource utilization when compared to highly tuned matrix-multiplication (GEMM) kernels—reporting about 35\% compared to the 80–90\% utilization attained on devices like the Hopper H100. One likely reason relates to implementation differences, particularly the continued use of Ampere-optimized instructions for Tensor Cores rather than those introduced with the Hopper architecture. Prior works—including ThunkerKitten~\citep{spector2024thunder} and cuDNN 9~\citep{cudnn9}—have demonstrated that attention computations can be accelerated and implementations streamlined when capitalizing on Hopper-specific instruction sets and tiling-based techniques.

At its core, the algorithm underlying \textsc{FlashAttention-2} operates within a streamlined synchronous paradigm, lacking deliberate integration of asynchrony or low-precision computation in its framework. Asynchrony emerges due to advances in hardware specialization that boost the efficiency of central tasks in machine learning workloads: for instance, matrix multiplication is delegated to dedicated hardware (Tensor Cores), and memory operations to specific modules (Tensor Memory Accelerator -- TMA), thereby decoupling these activities from the main CUDA cores, which handle control, integer, and conventional floating point operations. Meanwhile, reduced numerical precision—such as FP8 in Hopper and FP4 in Blackwell, building upon earlier moves to FP16 (with Pascal, 2017) and BF16 (in Ampere, 2020)—has repeatedly proven its effectiveness in increasing computational throughput (by factors of two or four) within fixed power and silicon constraints. We further outline the resources provided by Hopper in these regards in \cref{subsec:hardware}.

A significant technical hurdle lies in revamping \textsc{FlashAttention-2} to leverage these hardware advances. Incorporating asynchrony entails orchestrating concurrent progress of matrix multiplication and softmax even in the presence of data dependencies between them, while harnessing low-precision necessitates precise handling of quantization, especially for rare features in LLMs that can otherwise introduce substantial error~\citep{dettmers2208llm, sun2024massive}.

Addressing this, we introduce \textsc{FlashAttention-3}, which integrates and advances three novel strategies aimed at maximizing performance on the latest GPU architectures:\footnote{Our primary exposition is based on NVIDIA’s Hopper architecture, though the algorithm extends to any platform that supports comprehensive asynchronous operations and flexible low-precision data types.}

\begin{enumerate}

\item \textbf{Producer-Consumer asynchrony:} We introduce a software pipelining technique, specialized at the warp level, which leverages the asynchronous nature of both data transfers and Tensor Core computation. By assigning the roles of data producers and consumers to different warps, the approach enhances the capability of the algorithm to obscure the latency associated with memory accesses and instruction dispatching.
\item \textbf{Hiding softmax under asynchronous block-wise GEMMs:} Our method enables the overlap of low-throughput softmax-related computations—such as floating-point fused multiply-add and exponential operations—with asynchronously issued WGMMA GEMM instructions. This required modifications to the \textsc{FlashAttention-2} algorithm to eliminate some of the serial dependencies between softmax operations and GEMMs. In particular, for the two-stage algorithm variant, while softmax is being computed on a given block of the score matrix, an asynchronous WGMMA proxy processes the subsequent block.
\item \textbf{Hardware-accelerated low-precision GEMM:} We redesign the forward pass to fully utilize FP8 Tensor Cores for GEMM computations, which almost doubles the sustained TFLOPs/s in practice. Achieving this involves addressing distinct WGMMA layout constraints for storing FP32 accumulators and FP8 operand matrices in memory. Block quantization and incoherent processing are employed to counterbalance potential accuracy degradation associated with lowering to FP8 precision.
\end{enumerate}

To empirically assess the effectiveness of our approach, we conduct benchmarks of \textsc{FlashAttention-3} on the H100 SXM5 GPU using a diverse set of parameter configurations. Our evaluation demonstrates the following: (1) In the forward computation, FP16 delivers a 1.5-2.0$\times$ acceleration over \textsc{FlashAttention-2} (with peak performance up to 740 TFLOPs/s), and yields a 1.5-1.75$\times$ speedup in the backward computation; (2) FP8 reaches nearly 1.2 PFLOPs/s; (3) For long sequences, FP16 consistently outperforms, while FP8 remains highly competitive\footnote{Specifically, with head dimension 64, FP8 in \textsc{FlashAttention-3} surpasses other variants, whereas for head dimensions 128 and 256 it matches performance in cases without causal masking but lags behind when causal masking is present.} with the leading attention mechanism available in NVIDIA's cuDNN library. We further verify that FP16 in \textsc{FlashAttention-3} achieves the same level of numerical precision as \textsc{FlashAttention-2}, and it outperforms conventional attention routines, owing to the preservation of intermediate values (such as the rescaled softmax outputs) in FP32. Additionally, FP8 with block quantization and incoherent processing in \textsc{FlashAttention-3} attains 2.6$\times$ improved accuracy compared to standard attention utilizing per-tensor quantization, particularly in contexts where outlier features are present.

We have publicly released \textsc{FlashAttention-3} under a liberal license\footnote{\textsc{FlashAttention-3} is publicly accessible at \texttt{https://github.com/Dao-AILab/flash-attention}} and intend to integrate it with both PyTorch and Hugging Face ecosystems, aiming to benefit a broad community of practitioners and researchers.

\section{Background: Multi-Head Attention and GPU Characteristics}
\label{sec:background}

\subsection{Multi-Head Attention}
\label{subsec:multi_head_attn}

Consider $\mathbf{Q}, \mathbf{K}, \mathbf{V} \in \mathbb{R}^{N \times d}$ as the query, key, and value matrices provided to an individual attention head, with $N$ denoting the sequence length and $d$ the dimension per head. The resulting attention output $\mathbf{O}$ is obtained via:

\begin{equation*}
  \mathbf{S} = \alpha \mathbf{Q} \mathbf{K}^\top \in \mathbb{R}^{N \times N}, \quad \mathbf{P} = \mathrm{softmax}(\mathbf{S}) \in \mathbb{R}^{N \times N}, \quad \mathbf{O} = \mathbf{P}\mathbf{V} \in \mathbb{R}^{N \times d},
\end{equation*}

Here, the $\mathrm{softmax}$ function is implemented along the rows, and the scaling parameter $\alpha$ is usually chosen as $1/\sqrt{d}$. To mitigate potential issues with numerical precision caused by the exponentiation, it is standard to subtract $\mathrm{rowmax}(\mathbf{S})$ from $\mathbf{S}$. In the case of multi-head attention (MHA), distinct projections for queries, keys, and values are learned for each head. This process is executed in parallel for all heads and batches, culminating in the aggregated output tensor.

Consider a scalar loss function $\phi$, and denote its gradient with respect to an argument as $\mathbf{d}(-) = \partial \phi / \partial (-)$. Given the upstream gradient $\mathbf{dO} \in \mathbb{R}^{N \times d}$, the gradients with respect to $\mathbf{Q}$, $\mathbf{K}$, and $\mathbf{V}$ can be derived using the chain rule:

\begin{align*}
  \mathbf{dV} &= \mathbf{P}^\top \mathbf{dO} \in \mathbb{R}^{N \times d} \\
  \mathbf{dP} &= \mathbf{dO} \mathbf{V}^\top \in \mathbb{R}^{N \times N} \\
  \mathbf{dS} &= \mathrm{dsoftmax} (\mathbf{dP}) \in \mathbb{R}^{N \times N} \\
  \mathbf{dQ} &= \alpha \mathbf{dS} \mathbf{K} \in \mathbb{R}^{N \times d} \\
  \mathbf{dK} &= \alpha \mathbf{dS}^\top \mathbf{Q} \in \mathbb{R}^{N \times d},
\end{align*}

In the formula above, the relation $\mathbf{d}s = (\mathrm{diag}(p) - p p^\top)\mathbf{d}p$ holds for $p = \mathrm{softmax}(s)$, with $\mathrm{dsoftmax}(\mathbf{dP})$ representing this computation applied to each row independently. Additionally, the described procedure is amenable to parallel execution over heads and batch elements during the MHA backward pass.

\subsection{GPU hardware characteristics and execution model}
\label{subsec:hardware}

We describe the aspects of the GPU's execution model relevant for \textsc{FlashAttention-3}, with a focus on the NVIDIA Hopper architecture as a concrete instantiation of this model.

\paragraph{Memory hierarchy:} The GPU architecture features a stratified arrangement of memory types, where the storage capacity is inversely proportional to the available bandwidth (\cref{tab:gpu-hierarchy})\footnote{\citet{luo2024benchmarking} cites a shared memory bandwidth of 128 bytes per clock cycle per SM; by aggregating this across 132 SMs and factoring in the boost clock rate of 1830 MHz, the total is derived.}. Global memory (GMEM), or HBM, refers to the off-chip DRAM that can be accessed by all streaming multiprocessors (SMs). Data housed in GMEM is automatically cached within an on-chip L2 cache layer. Beyond this, every SM is equipped with a relatively small, on-chip shared memory (SMEM), which is extensively banked and directly managed by the programmer. At the deepest level of locality, each SM possesses its own register file.

\paragraph{Thread hierarchy:} In the GPU programming paradigm, computational units are systematically organized into hierarchical groups referred to as threads. This hierarchy, progressing from the most granular to the broadest levels, includes individual threads, followed by warps (each containing 32 threads), warpgroups (consisting of 4 adjacent warps), threadblocks (also known as cooperative thread arrays or CTAs), threadblock clusters (introduced with Hopper), and finally grids.

The relationship between these two hierarchies is highly interconnected. Threads belonging to a single CTA are scheduled together on a single SM, and similarly, CTAs grouped within one cluster are scheduled together on the same GPC. Every thread within a CTA can directly access SMEM, while RMEM (with a maximum of 256 registers) remains exclusive to each individual thread.

\paragraph{Asynchrony and warp-specialization:}

GPUs achieve high throughput by leveraging parallelism and asynchronous operations, which effectively mask delays in memory access and instruction execution. In the Hopper architecture, the Tensor Memory Accelerator (TMA) is introduced as specialized hardware to facilitate asynchronous memory transfers between GMEM and SMEM \cite[\S7.29]{cuda}. Additionally, in contrast to previous architectures like Ampere, Hopper’s Tensor Core—accessible through the warpgroup-wide WGMMA instruction \cite[\S9.7.14]{ptx}—operates asynchronously and is capable of directly fetching input data from shared memory.

The inclusion of hardware mechanisms for asynchrony enables the implementation of warp-specialized kernels. In this design, warps within a CTA are partitioned into groups that exclusively perform either data transfer operations (producers) or computational tasks (consumers). This structural division enhances the compiler’s capacity to create highly optimized instruction schedules \citep{warp-specialization-2011}. Hopper further extends this flexibility by permitting registers to be reallocated dynamically among warpgroups using \verb|setmaxnreg| \cite[\S9.7.17.1]{ptx}. Consequently, warps dedicated to performing MMAs can be allotted a greater portion of RMEM compared to warps responsible only for TMA operations (which typically require participation from a single thread).

\paragraph{Low-precision number formats:}
\label{sec:low-precision-gpu}
Recent generations of GPUs incorporate specialized architectural features aimed at boosting performance for low-precision arithmetic. Notably, the WGMMA instruction is optimized to exploit FP8 Tensor Cores on Hopper architectures, enabling twice the throughput per SM relative to FP16 or BF16 operations.

Nonetheless, effective use of FP8 WGMMA requires careful attention to operand layout constraints. Consider a GEMM operation performing $A \times B^{\top}$, where $A$ is an $M\times K$ matrix and $B$ is an $N\times K$ matrix. Here, an operand (either $A$ or $B$) is described as \emph{mn-major} if memory is contiguous along the outer $M$ or $N$ dimension; conversely, it is labeled \emph{k-major} if the inner $K$-dimension is contiguous. For FP16 WGMMA instructions, input operands in SMEM can use either the mn-major or k-major layout. In contrast, FP8 WGMMA restricts operand input to the k-major arrangement only. Additionally, scenarios such as attention—where chained GEMMs are fused within one kernel—present challenges, as the required layouts for the FP32 accumulators and the FP8 operands can be incompatible, making successive FP8 WGMMA invocations difficult.

In the context of attention, these layout restrictions entail certain modifications to the design of an FP8 algorithm, which we describe in \cref{sec:algofp8}.

\subsection{Standard Attention and Flash Attention}

According to \citet{dao2022flashattention}, \textbf{standard attention} refers to a GPU-based attention computation in which the intermediate matrices $\mathbf{S}$ and $\mathbf{P}$ are explicitly stored in HBM. In contrast, the key innovation behind \textsc{FlashAttention} is the application of a local softmax reduction to eliminate costly intermediate memory accesses, effectively integrating the attention calculation into a single computational kernel. This localized softmax operation is implemented in lines \ref{code-ws:softmax_start}-\ref{code-ws:softmax_end} within the main consumer loop presented in \cref{alg:flash3_wgmma_ws_only}, along with the block-wise scaling applied to $\mathbf{O}$. A straightforward derivation verifying that this mechanism correctly computes $\mathbf{O}$ can be found in \cite[\S 2.3.1]{dao2023flashattention2}.

\section{FlashAttention-3: Algorithm}
\label{sec:algo}

This section presents an overview of the \textsc{FlashAttention-3} algorithm. We center our explanation on the forward pass, deferring details of the backward pass to~\cref{sec:algo_ws_bwd}. Initially, we discuss incorporating warp-specialization along with a circular SMEM buffer into the foundational framework of \textsc{FlashAttention-2}. Next, we illustrate how the asynchronous capabilities of WGMMA can be leveraged to establish a two-stage GEMM-softmax pipeline with overlapping computations. Lastly, we outline the necessary alterations for supporting FP8, addressing both layout compatibility and enhanced accuracy through block quantization and the use of incoherent processing.

\subsection{Producer-Consumer asynchrony through warp-specialization and pingpong scheduling}
\label{sec:algo_ws}

\paragraph{Warp-specialization}
Similar to \textsc{FlashAttention-2}, the forward computation in \textsc{FlashAttention-3} features extensive parallelism across batch size, the number of attention heads, and the length of the query sequence. Therefore, we describe the algorithm at the CTA granularity, where each CTA processes a tile $\mathbf{Q}_i$ of the query matrix and calculates the associated output tile $\mathbf{O}_i$. For clarity, we introduce the warp-specialization approach utilizing a circular SMEM buffer, temporarily omitting consideration of GEMM-softmax overlap. Denote $d$ as the dimension of each head and $N$ as the total sequence length. Selecting a fixed block size $B_r$ for the queries, we partition $\mathbf{Q}$ into $T_r = \lceil \frac{N}{B_r} \rceil$ blocks, labelled $\mathbf{Q}_1, ..., \mathbf{Q}_{T_r}$.

\begin{algorithm}[H]
    \caption{\small\label{alg:flash3_wgmma_ws_only}\textsc{FlashAttention-3} forward pass \textbf{excluding} intra-consumer overlap -- CTA-level perspective}
    \begin{algorithmic}[1]
\REQUIRE Matrices $\mathbf{Q}_i \in \mathbb{R}^{B_r \times d}$ and $\mathbf{K}, \mathbf{V} \in \mathbb{R}^{N \times d}$ are located in HBM, with key blocks sized $B_c$ where $T_c = \lceil \frac{N}{B_c} \rceil$.
\STATE Instantiate pipeline control for managing synchronization barriers, leveraging an $s$-stage circular SMEM buffering system.
\IF {current warpgroup is a producer}
\STATE Free the assigned set of registers.
\STATE Trigger transfer of $\mathbf{Q}_i$ from HBM into shared memory.
\STATE Upon data arrival, signal completion of $\mathbf{Q}_i$ load to the consumers.
\FOR{$0 \le j < T_c$}
    \STATE Await confirmation that the $(j\,\%\,s)$ buffer segment has been fully processed by consumers.
    \STATE Transfer $\mathbf{K}_j, \mathbf{V}_j$ from HBM into shared memory within the $(j\,\%\,s)$ buffer partition.
    \STATE When done, signal to consumers that the transfer of $\mathbf{K}_j, \mathbf{V}_j$ is finished.
\ENDFOR
\ELSE \STATE Allocate back a set of registers according to the consumer warp count.
\STATE Locally, set $\mathbf{O}_i = (0) \in \mathbb{R}^{B_r \times d}$, initialize $\ell_i = (0) \in \mathbb{R}^{B_r}$ and $m_i = (-\infty) \in \mathbb{R}^{B_r}$.
\STATE Pause execution until $\mathbf{Q}_i$ is available in shared memory.
\FOR{$0 \le j < T_c$}
\STATE Wait until $\mathbf{K}_j$ is accessible in shared memory.
\STATE Calculate $\mathbf{S}_i^{(j)} = \mathbf{Q}_i \mathbf{K}_j^T$ using SS-GEMM; execute commit and await result. \label{alg:ws_only_gemm1}
\STATE Cache $m_i^{\mathrm{old}} = m_i$ and update $m_i = \mathrm{max}(m_i^{\mathrm{old}}, \mathrm{rowmax}(\mathbf{S}_i^{(j)}))$. \label{code-ws:softmax_start}
\STATE Compute $\widetilde{\mathbf{P}}_i^{(j)} = \mathrm{exp}(\mathbf{S}_i^{(j)} - m_i)$, then update $\ell_i = \mathrm{exp}(m_i^{\mathrm{old}} - m_i)\,\ell_i + \mathrm{rowsum}(\widetilde{\mathbf{P}}_i^{(j)})$. \label{code-ws:softmax_end}
\STATE Wait until $\mathbf{V}_j$ is present in shared memory.
\STATE Update $\mathbf{O}_i = \mathrm{diag}(\mathrm{exp}(m_i^{\mathrm{old}} - m_i))^{-1}\mathbf{O}_i + \widetilde{\mathbf{P}}_i^{(j)} \mathbf{V}_j$ via RS-GEMM; execute commit and await result. \label{alg:ws_only_gemm2}
\STATE Indicate to the producer that the $(j\,\%\,s)$ buffer section is free.
\ENDFOR
\STATE Finalize $\mathbf{O}_i = \mathrm{diag}(\ell_i)^{-1} \mathbf{O}_i$ and calculate $L_i = m_i + \log(\ell_i)$.
\STATE Persist $\mathbf{O}_i$ and $L_i$ in HBM as the $i$th partition of output tensors $\mathbf{O}$ and $L$.
\ENDIF
\end{algorithmic}
\end{algorithm}

In our deployment of \cref{alg:flash3_wgmma_ws_only} on Hopper, we utilize \verb|setmaxnreg| to manage (de)allocation operations, employ TMA for transferring $\mathbf{Q}_i$ and $\{ \mathbf{K}_j, \mathbf{V}_j \}_{0 \leq j < T_c}$, and rely on WGMMA to perform the GEMM computations within the consumer mainloop. The SS or RS prefix designates whether shared memory or register file serves as the source for the first operand, respectively. When analyzing the control flow in \cref{alg:flash3_wgmma_ws_only}, it is important to observe that asynchronous TMA loads proceed independently without waiting for prior loads to finish. Additionally, during the initial $s$ iterations of the producer mainloop, the system refrains from issuing waits since the buffer is in the process of being populated.

\paragraph{Pingpong scheduling} Leveraging the asynchronous execution model of WGMMA and TMA in conjunction with warp specialization enables the potential to concurrently execute softmax operations for one warpgroup while another warpgroup performs GEMM. This scheduling approach is motivated by the observation that, on contemporary hardware accelerators, non-matmul computations are significantly less performant than their matmul counterparts. For instance, on the H100 SXM5 GPU, FP16 matmul achieves 989 TFLOPS, whereas special function units—such as those required for the exponential computations in softmax—reach only 3.9 TFLOPS.\footnote{According to the CUDA programming guide, each streaming multiprocessor (SM) can handle 16 special function operations per clock cycle; with 132 SMs running at 1830 MHz, this sums to 3.9 TFLOPS for special functions.} In the context of the FP16 attention forward pass with a head dimension of 128, the total FLOP count for matmul exceeds that for exponentials by a factor of 512, but the corresponding throughput for exponentials is 256 times lower. Consequently, exponential computation can occupy up to 50\% of the execution cycles compared to matmul. The disparity becomes even greater with FP8: while matmul throughput doubles, the throughput for exponentials remains unchanged.

The computation of the exponential function is offloaded to a dedicated hardware unit known as the multi-function unit. Ideally, we aim to have the exponential calculations occur concurrently while the Tensor Cores execute the matrix multiplication operations. To orchestrate this, we introduce synchronization barriers (\texttt{bar.sync} instructions), enforcing an execution order such that the GEMMs (specifically, GEMM1 for $\mathbf{P} \mathbf{V}$ in the current iteration, and GEMM0 for $\mathbf{Q} \mathbf{K}^\top$ in the subsequent iteration) assigned to warpgroup 1 are completed prior to those of warpgroup 2. Consequently, the softmax computation for warpgroup 1 takes place in parallel with the GEMMs in warpgroup 2. This setup alternates each round: warpgroup 2 handles softmax while warpgroup 1 reverts to GEMMs, effectively implementing a ``pingpong'' scheduling pattern. An example of this mechanism is depicted in~\cref{fig:pingpong_scheduling}. Although, in real workloads, this pingpong scheduling is not perfectly aligned as the schematic might suggest, our empirical measurements show notable performance gains (for instance, forward pass throughput in FP16 increases from 570 TFLOPS to the range of 620-640 TFLOPS when using a head dimension of 128 and a sequence length of 8192).

\paragraph{Attention variants} In the context of multi-query attention~\citep{shazeer2019fast} and grouped query attention~\citep{ainslie2023gqa}, we employ the strategy from \textsc{FlashAttention-2}, modifying the tensor access patterns to prevent unnecessary duplication of $\mathbf{K}$ and $\mathbf{V}$ tensors in HBM.

\subsection{Intra-warpgroup overlapping GEMMs and softmax}

Even within one warpgroup, we can overlap some instructions in the softmax with some instructions in the GEMMs. We describe one technique to do so.

Within the attention algorithm, sequential dependencies in the core loop restrict intra-iteration parallelism. Specifically, the (local) softmax computation (spanning lines \ref{code-ws:softmax_start} to \ref{code-ws:softmax_end}) depends on $\mathbf{S}_i^{(j)}$, which is generated by the initial GEMM operation. Subsequently, the second GEMM consumes the output $\widetilde{\mathbf{P}}_i^{(j)}$ from softmax as its input. Notably, the insertion of wait statements at lines \ref{alg:ws_only_gemm1} and \ref{alg:ws_only_gemm2} in \cref{alg:flash3_wgmma_ws_only} forces softmax and both GEMMs to execute in sequence. To mitigate this serialization, we can introduce register-based buffers and organize computations in a pipelined fashion across different iterations. Building upon this strategy, we introduce a two-stage\footnote{It is important to observe that while the overlap technique's staging may be upper-bounded by the circular SMEM buffer size $s$, it does not necessarily require $s$ stages.} GEMM-softmax pipeline algorithm as follows:

\begin{algorithm}[H]
    \caption{\small\label{alg:flash3_wgmma}\textsc{FlashAttention-3} consumer warpgroup forward pass}
    \begin{algorithmic}[1]
      \REQUIRE  Matrices $\mathbf{Q}_i \in \mathbb{R}^{B_r \times d}$ and $\mathbf{K}, \mathbf{V} \in \mathbb{R}^{N \times d}$ located in HBM, with a specified key block size $B_c$ and corresponding $T_c = \lceil \frac{N}{B_c} \rceil$.
      \STATE Dynamically reassign a set number of registers, adjusting according to the current count of consumer warps.
      \STATE Locally initialize on-chip variables: set $\mathbf{O}_i = (0) \in \mathbb{R}^{B_r \times d}$, $\ell_i = (0)$, and $m_i = (-\infty)$ in $\mathbb{R}^{B_r}$.
      \STATE Stall execution until both $\mathbf{Q}_i$ and $\mathbf{K}_0$ are available in shared memory.
      \STATE Utilize WGMMA to calculate $\mathbf{S}_{\mathrm{cur}} = \mathbf{Q}_i \mathbf{K}_0^T$. Ensure completion before proceeding.
      \STATE Free the initial buffer segment corresponding to $\mathbf{K}$.
      \STATE Using $\mathbf{S}_{\mathrm{cur}}$, derive $m_{i}$, $\tilde{\mathbf{P}}_{\mathrm{cur}}$, and $\ell_{i}$, then adjust the scale of $\mathbf{O}_i$ accordingly.
      \FOR{$1 \le j < T_c - 1$}
        \STATE Wait for the transfer of $\mathbf{K}_j$ into shared memory to finish.
        \label{code:mainloop_start}
        \STATE Compute $\mathbf{S}_{\mathrm{next}} = \mathbf{Q}_i \mathbf{K}_{j}^T$ with WGMMA. Initiate but do not synchronize at this point. \label{code:first_wgmma}
        \STATE Ensure that $\mathbf{V}_{j-1}$ is completely moved to shared memory before continuing.
        \STATE Execute $\mathbf{O}_{i} = \mathbf{O}_{i} + \tilde{\mathbf{P}}_{\mathrm{cur}} \mathbf{V}_{j-1}$ using WGMMA, without blocking for completion. \label{code:second_wgmma}
        \STATE Synchronize to complete the $\mathbf{Q}_i \mathbf{K}_{j}^T$ WGMMA step.
        \STATE With $\mathbf{S}_{\mathrm{next}}$, update $m_{i}$, calculate $\tilde{\mathbf{P}}_{\mathrm{next}}$, and refresh $\ell_{i}$. \label{code:softmax}
        \STATE Once the WGMMA operation for $\tilde{\mathbf{P}}_{\mathrm{cur}} \mathbf{V}_{j-1}$ is done, apply the appropriate scaling to $\mathbf{O}_i$.
        \STATE Free the stage $(j\,\%\,s)$ in the buffer for $\mathbf{K}$, as well as $(j-1\,\%\,s)$ for $\mathbf{V}$.
        \STATE Move $\mathbf{S}_{\mathrm{next}}$ content to $\mathbf{S}_{\mathrm{cur}}$ for the following iteration.
        \label{code:mainloop_end}
      \ENDFOR \STATE Stall until $\mathbf{V}_{T_c - 1}$ has been loaded into shared memory.
      \STATE Finalize via WGMMA: calculate $\mathbf{O}_{i} = \mathbf{O}_{i} + \tilde{\mathbf{P}}_{\mathrm{last}} \mathbf{V}_{T_c - 1}$, then ensure completion.

\STATE Epilogue: Adjust $\mathbf{O}_{i}$ according to $m_{i}$. Derive $L_{i}$ utilizing both $m_{i}$ and $\ell_{i}$. Store $\mathbf{O}_{i}$ and $L_{i}$ in HBM, filling the $i$-th position of $\mathbf{O}$ and $L$.

\end{algorithmic}
\end{algorithm}

\cref{alg:flash3_wgmma} substitutes the consumer segment of \cref{alg:flash3_wgmma_ws_only}, thereby enabling the full implementation of the \textsc{FlashAttention-3} procedure for FP16 formats. At an abstract level, we use WGMMA to represent asynchronous GEMM. In the main computation loop (spanning lines \ref{code:mainloop_start} through \ref{code:mainloop_end}), the execution of the second WGMMA call in iteration $j$ (line \ref{code:second_wgmma}) coincides with the evaluation of softmax for the subsequent iteration $j+1$ (line \ref{code:softmax}).

Although the pipelined structure presented earlier provides potential gains in performance, several real-world factors must be addressed:

\paragraph{Compiler reordering} The pseudocode illustrates an optimalized, straightforward execution sequence; however, in reality, the NVCC compiler often reorders instructions to optimize for hardware. This reordering can interfere with the intended interleaving of WGMMA and non-WGMMA instructions, which might reduce or negate the anticipated benefits of the designed pipeline. Examination of the resulting SASS code reveals that overlapping of operations does occur as intended (see Section~\ref{sec:2-stage-sass}).

\paragraph{Register pressure} Achieving high performance depends on avoiding register spilling, as spills introduce significant overhead. The two-stage pipelined approach, however, necessitates the use of additional registers: it requires reserving an extra $\mathbf{S}_{\mathrm{next}}$ variable, increasing register usage by $B_r \times B_c \times \text{sizeof}(\text{float})$ for every threadblock. Such increased register consumption can conflict with strategies such as increasing threadblock size, which also increases register demand. Profiling data should inform any decision regarding these competing resource needs.

\paragraph{3-stage pipelining} By building on the two-stage algorithm, we can envision a 3-stage pipeline, enabling overlap of the second WGMMA with the softmax operation. While this configuration might enhance Tensor Core usage further, it substantially increases register requirements, since another set of intermediate values must be retained. The resultant tension between maximizing tile sizes and deepening the pipeline becomes even more pronounced. An in-depth discussion of the 3-stage algorithm and empirical results is provided in ~\cref{sec:3-stage}.

\subsection{Low-precision with FP8}
\label{sec:algofp8}

\textbf{Efficiency: layout transformations.} Performing the forward computation for \textsc{FlashAttention-3} in FP8 precision introduces unique difficulties regarding tensor layout compliance that are not present for FP16.

Typically, the input tensors $\mathbf{Q}$, $\mathbf{K}$, and $\mathbf{V}$ are stored with contiguous memory in the head dimension. However, to meet the k-major requirement imposed by FP8 WGMMA on the second GEMM, it is necessary for $\mathbf{V}$—specifically, the tiles of $\mathbf{V}$ transferred into SMEM—to be contiguous along the sequence length dimension. Because the TMA load operation does not alter which dimension is contiguous in memory, there are two principal strategies: (1) performing an explicit transpose of $\mathbf{V}$ in global memory (GMEM) before further computation, or (2) carrying out a transpose of the required $\mathbf{V}$ tiles within the kernel itself after they have been moved into SMEM. For the first strategy, we can (1a) integrate the transpose with the epilogue of an earlier step, such as rotary embedding, or (1b) invoke a dedicated pre-processing kernel\footnote{A highly optimized transpose kernel can achieve device-level bandwidth efficiency~\citep{colfax_cutlass_transpose_2024}.} that swaps the strides for the head and sequence length dimensions. However, (1a) presents integration challenges for general-purpose libraries, while (1b) proves inefficient and introduces excessive overhead for memory-bound scenarios like inference.

In the case of FP8 \textsc{FlashAttention-3}, we select option (2) as our approach. To accomplish the in-kernel transpose, we utilize the LDSM (\verb|ldmatrix|) and STSM (\verb|stmatrix|) instructions. These enable a warp of threads to cooperate in transferring data between SMEM and RMEM in blocks of 128 bytes.\footnote{The PTX specification identifies LDSM/STSM as instructions for copying $8 \times 8$ matrices composed of 16-bit elements \cite[\S 9.7.13.4.15-16]{ptx}. By pairing 8-bit elements, however, LDSM/STSM can be applied in the FP8 context. That said, the transpose variants of these instructions are unable to separate packed 8-bit values, thus necessitating additional register-level manipulations between LDSM and STSM for tile-wise transposes; we omit the low-level specifics here.} Both LDSM and STSM efficiently use registers, permitting execution within the producer warpgroup, and additionally provide layout transposition during memory copy operations. Further, starting from the second iteration, the transpose for the upcoming $\mathbf{V}$ tile can be performed concurrently—hidden behind—the two WGMMAs operating on the prior $\mathbf{V}$ and present $\mathbf{K}$ tiles.

Second, we note that, in contrast to FP16, the arrangement in memory of the FP32 accumulator used by an FP8 WGMMA diverges from the organization assumed for operand A when it is resident in registers. Illustrations of both layouts are provided in \cref{fig:rmem_accum} and \cref{fig:rmem_operand}, which show the sequence in which register entries are distributed across threads. Employing byte permute instructions allows us to reformat the accumulator from the initial WGMMA into the organization required for subsequent WGMMA computations, aligning it with the structure of the $\mathbf{V}$ tile as generated by the in-kernel transpose. More precisely, as indicated in \cref{fig:rmem_accum}, this involves rearranging the bytes in the order
$$\{ \verb|d0 d1 d4 d5 d2 d3 d6 d7| \},$$
whereby this register permutation is repeated every 8 bytes. With regard to the logical configuration of the $\mathbf{P}$ tile, this operation induces a permutation of its columns (so, for example, columns $0189$ are mapped to the first four columns). To ensure WGMMA produces the correct output tile, the in-kernel transpose can be programmed to yield a corresponding row reordering of the $\mathbf{V}$ tile. \footnote{Leveraging this additional flexibility provided by performing the transpose inside the kernel dispenses with the need for shuffle instructions to reassign registers between threads, a practice previously detailed in~\citep{colfax_fp8_flashattention_2024}.}

\textbf{Accuracy: block quantization and incoherent processing.} The FP8 (e4m3) data type dedicates 3 bits to the mantissa and 4 bits to the exponent, resulting in increased quantization error relative to FP16 or BF16. The challenge is exacerbated in large-scale models due to the frequent occurrence of extreme outlier values~\citep{dettmers2208llm, sun2024massive}, whose magnitudes substantially exceed most entries and complicate quantization strategies. Commonly, per-tensor scaling~\citep{micikevicius2022fp8} is adopted, meaning a single scaling factor is applied per tensor—one each for $\mathbf{Q}$, $\mathbf{K}$, and $\mathbf{V}$. To mitigate accuracy loss in attention computations when operating in FP8, we leverage two main approaches:

\begin{enumerate}

\item \textbf{Block quantization}: In this approach, a single scaling factor is maintained per block. Specifically, for each of $\mathbf{Q}$, $\mathbf{K}$, and $\mathbf{V}$, the tensors are partitioned into blocks of shape either $B_r \times d$ or $B_c \times d$, with quantization performed independently on each block. This quantization step can be seamlessly combined with operations immediately preceding the attention calculation (such as rotary embedding) without introducing additional latency, due to the memory bandwidth limitations of rotary embedding. Notably, as \textsc{FlashAttention-3} natively processes data in blocks, each block in $\mathbf{S}$ can be rescaled post-quantization, effectively incorporating the block quantization without incurring further computational overhead.

\item \textbf{Incoherent processing}: To mitigate the effects of outliers, $\mathbf{Q}$ and $\mathbf{K}$ are pre-multiplied by a random orthogonal matrix $\mathbf{M}$ prior to FP8 quantization. The orthogonality property ($\mathbf{M} \mathbf{M}^\top = I$) ensures that $(\mathbf{Q} \mathbf{M}) (\mathbf{K} \mathbf{M})^\top$ is equivalent to $\mathbf{Q} \mathbf{K}^\top$, preserving the integrity of the attention computation. This procedure disperses outlier values because the entries of $\mathbf{Q} \mathbf{M}$ and $\mathbf{K} \mathbf{M}$ become random linear combinations of the original entries, thereby lowering the potential quantization error. In practice, following \citet{chee2024quip} and \citet{tseng2024quip}, $\mathbf{M}$ is constructed as the product of random diagonal matrices with $\pm 1$ entries and a Hadamard matrix. This form allows for multiplication in $O(d \log d)$ complexity, as opposed to $O(d^2)$, and can similarly be combined with the rotary embedding with no additional computational expense.

We validate that these two techniques reduces numerical error by up to 2.6$\times$ in \cref{sec:numerical_error}.

\section{Empirical Validation}
\label{sec:exp}

We use the primitives from CUTLASS~\citep{Thakkar_CUTLASS_2023} such as WGMMA and TMA abstractions to implement \textsc{FlashAttention-3} and evaluate its efficiency and accuracy.

\begin{itemize}

\item \textbf{Benchmarking attention.} We evaluate the performance of \textsc{FlashAttention-3} by timing its execution for various sequence lengths, comparing these results to those from the standard PyTorch implementation, \textsc{FlashAttention-2}, the Triton-optimized version of \textsc{FlashAttention-2} (leveraging H100-specific instructions), and a vendor-optimized \textsc{FlashAttention-2} leveraging cuDNN for H100 GPUs. Our results demonstrate that \textsc{FlashAttention-3} achieves speedups of up to 2.0$\times$ relative to \textsc{FlashAttention-2}, and 1.5$\times$ compared to \textsc{FlashAttention-2} in Triton. Furthermore, \textsc{FlashAttention-3} attains throughput as high as 740 TFLOPs/s, corresponding to 75\% of the theoretical TFLOPs/s peak on H100 GPUs.

\item \textbf{Ablation study.} Through an ablation analysis, we establish that optimizations such as warp-specialization and the integration of GEMM-softmax pipelining are significant contributors to the observed acceleration in \textsc{FlashAttention-3}.

\item \textbf{Accuracy of FP8 attention.} We verify that employing block quantization in conjunction with incoherent processing effectively decreases numerical errors in FP8 \textsc{FlashAttention-3} computations by a factor of 2.6$\times$.

\subsection{Benchmarking Attention}
\label{subsec:benchmark_attn}

We evaluate the execution time of various attention algorithms on an H100 80GB SXM5 GPU across multiple configurations, specifically toggling the causal mask and using head dimensions of either 64 or 128, all with FP16 input precision. The outcomes, presented in~\cref{fig:benchmark_attn_fwd} and~\cref{fig:benchmark_attn_bwd}, reveal that \textsc{FlashAttention-3} achieves forward pass speeds roughly 1.5-2.0$\times$ greater than those of \textsc{FlashAttention-2}, and accelerates the backward pass by approximately 1.5-1.75$\times$. When benchmarked against a conventional attention implementation, \textsc{FlashAttention-3} offers a speedup ranging from 3$\times$ up to 16$\times$. Notably, for sequence lengths at or above 1,000 tokens, \textsc{FlashAttention-3} outperforms even highly optimized proprietary solutions, such as the vendor’s cuDNN library for the H100 architecture.

\paragraph{Benchmark settings:} Sequence lengths are adjusted among 512, 1k, up to 16k, while maintaining a total of 16k tokens per batch by configuring the batch size accordingly. The hidden dimension remains fixed at 2048. The head dimension is set to 64, 128, or 256, corresponding to 32, 16, or 8 heads, respectively. FLOPs for the forward pass are computed using:

\begin{equation*}
  4 \cdot \text{seqlen}^2 \cdot \text{head dimension} \cdot \text{number of heads}.
\end{equation*}

When applying causal masking, we halve this value to reflect that roughly only 50\% of the elements are computed. To determine the FLOPs for the backward pass, we scale the forward pass FLOPs by a factor of 2.5, owing to the presence of 2 matrix multiplications during the forward computation and 5 during the backward computation (because of recomputation).

We also evaluate the forward pass runtime using FP8 under analogous experimental conditions. The outcomes for a headdim of 256 are depicted in ~\cref{fig:benchmark_attn_fp8}, while comprehensive results are provided in \cref{sec:benchmark_attn_fp8_full}.

\subsection{Ablation Study: 2-Stage Pipelining Experiments}
\label{sec:2-stage-pipelining-experiments}

We conduct an ablation study on both the 2-stage WGMMA-softmax pipelining and warp-specialization for the non-causal FP16 \textsc{FlashAttention-3}, using fixed parameters $\{\text{batch}, \text{seqlen}, \text{nheads}, \text{hdim}\} = \{ 4, 8448, 16, 128\}$. The findings, as shown in~\cref{table:ablation_pipelining}, indicate that our algorithmic enhancements—specifically asynchrony through warp-specialization and the concurrency of GEMM with softmax computations—yield substantial performance gains, increasing throughput from 570 TFLOPs to 661 TFLOPs.

\subsection{Numerical Error Validation}
\label{sec:numerical_error}

Given ongoing discussions regarding numerical error~\citep{golden2024flash} in
\textsc{FlashAttention}, we evaluate the accuracy of \textsc{FlashAttention-2}, \textsc{FlashAttention-3}, and a conventional attention implementation relative to a reference implementation using FP64 precision. To mimic the presence of extreme values and activations commonly found in LLMs~\citep{dettmers2208llm,
  sun2024massive}, we sample the elements of $\mathbf{Q}, \mathbf{K}, \mathbf{V}$ from the following distribution:

\begin{equation*}
  \mathcal{N}(0, 1) + \mathcal{N}(0, 100) \cdot \mathrm{Bernoulli}(0.001).
\end{equation*}

Specifically, every matrix entry follows a normal distribution with a mean of zero and a standard deviation of 1; however, for 0.1\% of the elements, we additionally incorporate an independent Gaussian component with a standard deviation of 10. The root mean squared error (RMSE) is reported in~\cref{table:numerical_error}. When running in FP16, both \textsc{FlashAttention-2} and \textsc{FlashAttention-3} yield RMSE values that are 1.7$\times$ smaller relative to the standard baseline because the intermediate values (i.e., the softmax output) are retained in FP32 precision. In contrast, the baseline FP8 attention employs per-tensor scaling and executes matrix multiplications with an FP32 accumulator while softmax intermediates are stored in FP16. Owing to its use of block quantization and incoherent computation, \textsc{FlashAttention-3} operating in FP8 delivers a 2.6$\times$ improvement in accuracy compared to the aforementioned baseline.

\section{Dicussion, Limitations, Conclusion}
\label{sec:discussion}

Through the development of \textsc{FlashAttention-3}, we have shown that leveraging modern programming paradigms and exploiting hardware capabilities—namely asynchrony and low-precision computation—can substantially enhance both the performance and fidelity of attention mechanisms. Our approach achieves a 1.5-2.0$\times$ acceleration over \textsc{FlashAttention-2} and diminishes FP8 numerical errors by 2.6$\times$ relative to traditional per-tensor quantization methods. Nevertheless, there remain several directions for future improvement: tailoring optimization for LLM inference tasks, adopting a persistent kernel design for the FP8 implementation,\footnote{In our evaluations, the FP16 variant of \textsc{FlashAttention-3} utilizes both a persistent kernel and an advanced load balancing approach, whereas these features are absent in the FP8 version. This discrepancy partly accounts for FP8 \textsc{FlashAttention-3} underperforming on tasks with small sequence lengths and causal masks compared to the FP8 cuDNN implementations.} and rigorously characterizing the implications of low-precision attention within large-scale model training. While this study centers on Hopper GPUs, the methodological advances introduced are expected to generalize to alternative hardware accelerators. We anticipate that making attention both faster and more accurate will open up new opportunities for long-context applications.

\end{document}